\documentclass{article}
\usepackage[UTF8]{ctex}
\usepackage{amsmath,amssymb,amsfonts}
\usepackage{enumitem}
\usepackage{geometry}
\geometry{a4paper, margin=1in}
% === 错题本排版优化（不参与题目标识） ===
\setlength{\parindent}{0pt}
\setlength{\parskip}{0.35em plus 0.10em minus 0.08em}
\setlength{\emergencystretch}{3em}
\setlist{itemsep=0.25em, topsep=0.40em, parsep=0pt}
\allowdisplaybreaks[2]
% === 错题本排版优化结束 ===
\begin{document}
\section*{吉林大学2026年数学分析真题}
\section*{77．吉林大学2026年数学分析试卷}

\subsection{题目1}
1． \(\lim _{n \rightarrow \infty} n^{3}\left[\sin \left(\tan \frac{1}{n}\right)-\sin \frac{1}{n}\right]\) ．

\subsection{题目2}
2． \(\lim _{n \rightarrow \infty} \frac{\sqrt[n]{(3 n)!}}{n^{3}}\) ．

\subsection{题目3}
3． \(\lim _{n \rightarrow \infty}\left[\left(\frac{\ln ((n+1)!)}{n+1}-\frac{\ln (n!)}{n}\right) \sum_{k=1}^{n} \sqrt[k]{k^{2}+\sin k}\right]\) ．

\subsection{题目4}
4． \(\lim _{x \rightarrow 0} \frac{\left(\sqrt{1+x^{2}}-\sqrt{\cos x}\right) \tan x^{2}}{\int_{x^{2}}^{x} \ln \left(1+t^{3}\right) \mathrm{d} t}\) ．

\subsection{题目5}
5． \(\lim _{n \rightarrow \infty} n \int_{0}^{1} x^{n} \ln (1+x) \mathrm{d} x\) ．

\subsection{题目6}
6．设 \(u, v\) 是由 \(\left\{\begin{array}{l}u^{2}-v=3 x+y \\ u-2 v^{2}=x-2 y\end{array}\right.\) 确定的隐函数，求 \(\frac{\partial u}{\partial x}, \frac{\partial v}{\partial y}\) ．

\subsection{题目7}
7．计算 \(\iiint_{D}\left(x^{2}+y^{2}+z^{2}\right) \mathrm{d} x \mathrm{~d} y \mathrm{~d} z\) ，其中 \(D\) 是两个球 \(x^{2}+y^{2}+z^{2} \leq R^{2}, x^{2}+y^{2}+(z-R)^{2} \leq R^{2}\) 的公共部分。

\subsection{题目8}
8．计算第二型曲线积分 \(\int_{L} \frac{y \mathrm{~d} x+(1-x) \mathrm{d} y}{(x-1)^{2}+y^{2}}\) ，其中 \(L\) 为圆周 \(x^{2}+y^{2}=4\) ，方向取顺时针。

二．当 \(x, y, z\) 都大于 0 时，求 \(u=\ln x+2 \ln y+3 \ln z\) 在球面 \(x^{2}+y^{2}+z^{2}=6\) 上的最大值．

三、证明函数项级数 \(\sum_{n=1}^{\infty} \sin n x \ln \left(1+\frac{1}{n}\right)\) 在 \((0,2 \pi)\) 上逐点收敛但不一致收敛．

四．求数项级数 \(\sum_{n=1}^{\infty} \frac{1}{n(2 n-1)}\left(-\frac{1}{4}\right)^{n}\) 的和．

五。讨论参数 \(p\) 对广义积分 \(\int_{0}^{+\infty} \frac{\cos \left(x-\frac{1}{x}\right)}{x^{p}} \mathrm{~d} x\) 的敛散性影响，何时绝对收敛，何时条件收敛，何时发散？

六．设 \(f(x)\) 在 \([1,+\infty)\) 上有定义，且存在常数 \(L\) ，使得对任意 \(x_{1}, x_{2} \in[1,+\infty)\) ，都有

\[
\left|f\left(x_{1}\right)-f\left(x_{2}\right)\right| \leq L\left|x_{1}-x_{2}\right| .
\]

试证存在 \(X \in[1,+\infty)\) ，使得 \(\frac{f(x)}{x+\ln \left(1+\frac{1}{x}\right)}\) 在 \([X,+\infty)\) 上一致连续。

七．假设 \(f(x)\) 在 \([0,2]\) 上四次连续可微，设

\[
g(x)=\frac{f(0)}{2}(x-1)(x-2)-f(1) x(x-2)+\frac{f(2)}{2} x(x-1) .
\]

（1）证明：任给 \(x \in[0,2]\) ，都存在 \(\xi \in(0,2)\) ，使得

\[
f(x)-g(x)=\frac{f^{\prime \prime \prime}(\xi)}{6} x(x-1)(x-2) .
\]

（2）证明：存在 \(\eta \in(0,2)\) ，使得 \(\int_{0}^{2}[f(x)-g(x)] \mathrm{d} x=-\frac{1}{90} f^{\prime \prime \prime \prime}(\eta)\) ．\\
（3）定义

\[
\varepsilon_{n}=\int_{0}^{2} f(x) \mathrm{d} x-\frac{1}{3 n} \sum_{i=1}^{n}\left[f\left(\frac{2(i-1)}{n}\right)+4 f\left(\frac{2 i-1}{n}\right)+f\left(\frac{2 i}{n}\right)\right] .
\]

证明： \(\lim _{n \rightarrow \infty} n^{4} \varepsilon_{n}=-\frac{1}{180}\left[f^{\prime \prime \prime}(2)-f^{\prime \prime \prime}(0)\right]\) ．
\end{document}
